% In-class activity: scheduling three FreeRTOS tasks on the STM32 by hand.
% Self-contained -- no Python data generation needed (the two timing charts
% are drawn directly with TikZ). Compile normally for the student handout,
% or with \def\keymode{1} before \input (same pretex trick used by the
% other activities/homeworks in this repo, see Generate_Scheduler_Activity.py)
% to overlay the solved task bars as the instructor key.
\documentclass[11pt]{article}
\usepackage[landscape,margin=0.5in]{geometry}
\usepackage{enumitem}
\usepackage{xcolor}
\usepackage{fancyhdr}
\usepackage{tikz}
\usetikzlibrary{shapes.geometric,arrows.meta}

% Colors matched to the original slide deck (Google Drawings blue/orange),
% plus two more reused for the Task 3 and Idle bars in the instructor key.
\definecolor{taskblue}{RGB}{74,134,232}
\definecolor{taskorange}{RGB}{255,171,64}
\definecolor{taskgreen}{RGB}{52,168,83}
\definecolor{idlegray}{RGB}{189,189,189}

\tikzset{
  t1burst/.style={star, star points=10, star point ratio=0.55,
    fill=taskblue, text=white, draw=none, font=\bfseries\footnotesize,
    align=center, minimum size=1.35cm, inner sep=0pt},
  t2burst/.style={star, star points=10, star point ratio=0.55,
    fill=taskorange, text=black, draw=none, font=\bfseries\footnotesize,
    align=center, minimum size=1.35cm, inner sep=0pt},
}

% Solved-schedule bar for the instructor key: a filled rectangle from time
% #1 to time #2 (in 10 us steps, matching \schedAxis's x coordinates) in
% lane #3 (4=Task 1, 3=Task 2, 2=Task 3, 1=Idle), filled with color #4.
\newcommand{\taskbar}[4]{%
  \draw[fill=#4, draw=black!55] (#1,#3-0.32) rectangle (#2,#3+0.32);
}

\ifdefined\keymode
  \newcommand{\docsuffix}{ -- Instructor Key}
\else
  \newcommand{\docsuffix}{}
\fi

\setlength{\parindent}{0pt}
\setlist[itemize]{leftmargin=*, itemsep=2pt, topsep=2pt}
\pagestyle{fancy}
\fancyhf{}
\lhead{ECGR 4101/5101 -- Advanced Embedded Systems}
\rhead{In-Class Activity: Schedulers\docsuffix}
\renewcommand{\headrulewidth}{0.4pt}
\setlength{\headheight}{14pt}
\setlength{\headsep}{10pt}

% The time axis (0--230 us in 10 us steps) and the four task/idle lanes,
% shared by both scenarios. Gridlines run from the axis up to the row of
% event markers.
\newcommand{\schedAxis}{%
  \foreach \t in {0,10,...,230} {
    \draw[gray!55, dash pattern=on 2pt off 2pt] (\t/10,0) -- (\t/10,4.6);
    \node[below, font=\scriptsize, gray!70] at (\t/10,0) {\t};
  }
  \node[anchor=east, font=\scriptsize, gray!70] at (-0.2,-0.5) {Time (us):};
  \node[anchor=east, font=\large] at (-0.3,4) {Task 1:};
  \node[anchor=east, font=\large] at (-0.3,3) {Task 2:};
  \node[anchor=east, font=\large] at (-0.3,2) {Task 3:};
  \node[anchor=east, font=\large] at (-0.3,1) {Idle:};
}

\begin{document}

{\Large\bfseries In-Class Activity: Schedulers}\hfill
Name: \underline{\hspace{6cm}}\quad Date: \underline{\hspace{2.5cm}}

\bigskip
\textbf{Part 1 -- Preemption, No Time Slicing}

\smallskip
\textbf{Description}

\smallskip
You have three tasks that are being managed in threads by FreeRTOS on the
STM32. The scheduler is configured to use \textbf{no time slicing} and
\textbf{preemption}.

\smallskip
Each of these tasks have the following properties:
\begin{itemize}
  \item Task 1: A medium-level priority task. Runs for 20 us. Event-driven.
  \item Task 2: A high-priority task. Runs for 10 us. Event-driven.
  \item Task 3: A low priority task. Runs for 30 us. Continuous task that is
    called every 20 us after completion of its previous completion. Starts
    at time 0.
\end{itemize}

\bigskip
\begin{center}
\begin{tikzpicture}[x=0.95cm, y=1cm, >={Stealth[length=3mm]}]
  \schedAxis

  \ifdefined\keymode
    \taskbar{0}{3}{2}{taskgreen}   % Task 3 release 1: runs 0-30 us
    \taskbar{3}{5}{1}{idlegray}    % Idle 30-50 us
    \taskbar{5}{6}{2}{taskgreen}   % Task 3 release 2 starts @50, runs 50-60...
    \taskbar{6}{8}{4}{taskblue}    % ...preempted by Task 1 event @60, 60-80
    \taskbar{8}{10}{2}{taskgreen}  % ...Task 3 resumes, finishes @100
    \taskbar{10}{12}{1}{idlegray}  % Idle 100-120 us
    \taskbar{12}{14}{2}{taskgreen} % Task 3 release 3 starts @120, runs 120-140...
    \taskbar{14}{15}{4}{taskblue}  % ...preempted by Task 1 event @140, 140-150...
    \taskbar{15}{16}{3}{taskorange}% ...preempted by Task 2 event @150, 150-160
    \taskbar{16}{17}{4}{taskblue}  % ...Task 1 resumes its last 10 us, 160-170
    \taskbar{17}{18}{2}{taskgreen} % ...Task 3 resumes its last 10 us, finishes @180
    \taskbar{18}{20}{1}{idlegray}  % Idle 180-200 us
    \taskbar{20}{23}{2}{taskgreen} % Task 3 release 4 starts @200, runs to 230
  \fi

  \node[t1burst] (p1ev1) at (6,5.6) {Task 1\\Event!};
  \draw[->, taskblue, line width=2.2pt] (p1ev1.south) -- (6,4.6);

  \node[t1burst] (p1ev2) at (14,5.4) {Task 1\\Event!};
  \draw[->, taskblue, line width=2.2pt] (p1ev2.south) -- (14,4.6);

  \node[t2burst] (p1ev3) at (15,6.9) {Task 2\\Event!};
  \draw[->, taskorange, line width=2.2pt] (p1ev3.south) -- (15,4.6);
\end{tikzpicture}
\end{center}

\newpage
\textbf{Part 2 -- Preemption with Time Slicing}

\smallskip
\textbf{Description}

\smallskip
You have three tasks that are being managed in threads by FreeRTOS on the
STM32. The scheduler is configured to use \textbf{time slicing} and
\textbf{preemption}. Time slices are configured to be 20 us long.

\smallskip
Each of these tasks have the following properties:
\begin{itemize}
  \item Task 1: A medium-level priority task. Runs for 40 us. Event-driven.
  \item Task 2: A high-priority task. Runs for 10 us. Event-driven.
  \item Task 3: A medium-priority task. Runs for 60 us. Continuous task that
    is called every 20 us after completion of its previous completion.
    Starts at time 0.
\end{itemize}

\bigskip
\begin{center}
\begin{tikzpicture}[x=0.95cm, y=1cm, >={Stealth[length=3mm]}]
  \schedAxis

  \ifdefined\keymode
    \taskbar{0}{4}{2}{taskgreen}   % Task 3 release 1 runs alone, 0-40 us
    \taskbar{4}{6}{4}{taskblue}    % Task 1 ready @30, round-robins in @40
    \taskbar{6}{8}{2}{taskgreen}   % back to Task 3's turn @60, finishes @80
    \taskbar{8}{10}{4}{taskblue}   % Task 1 alone again, finishes @100
    \taskbar{10}{11}{2}{taskgreen} % Task 3 release 2 starts @100, runs 100-110...
    \taskbar{11}{12}{3}{taskorange}% ...preempted by Task 2 event @110, 110-120
    \taskbar{12}{17}{2}{taskgreen} % ...Task 3 resumes alone, finishes @170
    \taskbar{17}{18}{1}{idlegray}  % Idle 170-180 us
    \taskbar{18}{20}{4}{taskblue}  % Task 1 event @180, runs alone until Task 3...
    \taskbar{20}{22}{2}{taskgreen} % ...release 3 (ready @190) round-robins in @200
    \taskbar{22}{23}{4}{taskblue}  % back to Task 1's turn @220 (continues past 230)
  \fi

  \node[t1burst] (p2ev1) at (3,5.6) {Task 1\\Event!};
  \draw[->, taskblue, line width=2.2pt] (p2ev1.south) -- (3,4.6);

  \node[t2burst] (p2ev2) at (11,5.6) {Task 2\\Event!};
  \draw[->, taskorange, line width=2.2pt] (p2ev2.south) -- (11,4.6);

  \node[t1burst] (p2ev3) at (18,5.6) {Task 1\\Event!};
  \draw[->, taskblue, line width=2.2pt] (p2ev3.south) -- (18,4.6);
\end{tikzpicture}
\end{center}

\end{document}
